\subsection{Upper idempotents detect the necessary coefficients}
The two-sided context construction admits the stronger testing family
in Theorem~\ref{thm:upper-idempotent-quotient}. The lifting step remains
the ordinary nilpotent correction already proved above and attributed
to \cite[Lemma~10.32.7, Tag~05BU]{Stacks}.

\begin{proof}[Proof of Theorem~\ref{thm:upper-idempotent-quotient}]
For $e\in U(Q)$, the coefficient product lands in $F(e^2)=F(e)$.
Associativity and biadditivity make this a ring. Naturality applied to
$1\leq e$ shows that the transitioned unit acts as a two-sided
identity.

Each defining test for $K^{\mathrm{id}}_q$ is an additive homomorphism.
If $q\leq q'$, a test at $q'$ with target $e$ is the same test at $q$
followed by transport, since $pqr\leq pq'r\leq e$. Thus transitions
preserve the subgroups. Multiplication on either side preserves them
because a test of $axb$ can be regrouped as
\[
 u(axb)v=(ua)x(bv),
\]
with the same ordered support bound. The groups therefore form a
coefficient ideal, without a commutativity assumption. The unit test
at $e\in U(Q)$ gives $K^{\mathrm{id}}_e=0$. Conversely, every
coefficient ideal vanishing at these supports sends each displayed
context into its zero group at $e$, so it is contained in
$K^{\mathrm{id}}$. This proves maximality and the existence of the
quotient. The inverse image of a quotient ideal vanishing at every
$e\in U(Q)$ has the same property upstairs, since these fibers were
unchanged. Maximality makes that quotient ideal zero.

Fix $E^2=E$. For a vertex $v$, put $a=E_{vv}$. The diagonal
idempotence equation gives $a^2\leq a$, and consequently
\[
 e_v=a\vee1\in U(Q),\qquad
 e_v^2=a^2\vee a\vee1=a\vee1.
\]
For each edge support $b=E_{vw}$, the same matrix equation gives
\[
 e_vb=(E_{vv}\vee1)E_{vw}=E_{vw}=b.
\]
If $z\in K^{\mathrm{id}}_{E_{vv}}$, its transition to $F(e_v)$ is
zero. Naturality therefore shows that $zy$, transported to $F(E_{vw})$,
is zero for every $y\in F(E_{vw})$.

The additive matrix fiber $A_E$ is the possibly nonunital ring with
entries $F(E_{ij})$ and the prescribed matrix transports. Its kernel
$N_E$ has entries $K^{\mathrm{id}}_{E_{ij}}$. In a product of $n+1$
elements of $N_E$, expand an entry into contributions of paths
$i_0,\ldots,i_{n+1}$. Among the first $n+1$ vertices a vertex repeats,
so there is a nonempty cycle with a following edge. Transport the
cycle coefficient to $F(E_{vv})$, using $E^m=E$ for every positive
$m$. It remains in $K^{\mathrm{id}}_{E_{vv}}$. The preceding
annihilation identity kills that cycle and its following edge after
transport to the edge support. Multiplying the path prefix and suffix
and using naturality kills the original contribution. Thus
$N_E^{n+1}=0$.

For completeness, lift a quotient idempotent at its exact support
$E$ to $a\in A_E$. Its defect $d=a^2-a$ lies in $N_E$. Replacing
$a$ by $3a^2-2a^3$ changes the defect to $d^2(4d-3)$, so finitely
many repetitions produce an idempotent. The polynomial has no constant
term, hence is valid in $A_E$ even when this ring is nonunital.
To lift an isomorphism, take rectangular inverse intertwiners in the
quotient, lift their coefficients at the same supports, and cut them
by the two lifted idempotents $P,Q$. The resulting $X=PXQ$ and
$Y=QYP$ have composites $XY=P+\delta$, $YX=Q+\epsilon$, where the
errors lie in the respective nilpotent fiber ideals. Finite geometric
series invert both composites in their corners. Then
\[
 Y'=Y(XY)^{-1}=(YX)^{-1}Y
\]
satisfies $XY'=P$ and $Y'X=Q$. Existence and uniqueness of lifted
projective classes follow. Extension of scalars preserves the split
idempotent images and their direct sums, proving the assertions on
$V$ and $K_0$. All operations preserve exact support; zero amplitude
at a nonzero support has not become the global zero object.

The old upper-unit ideal vanishes at every upper idempotent, so
maximality includes it in $K^{\mathrm{id}}$. If $1$ is greatest,
$U(Q)=\{1\}$. The unit test and multiplication naturality give
$K^{\mathrm{id}}_q=\ker(F(q)\to F(1))$. As before, pair isomorphisms
transport all tests, but arbitrary-map functoriality has not been
asserted.
\end{proof}

\begin{example}[strictness and the scope of nilpotence]
\label{ex:upper-idempotent-strict}
Let $Q=\{\bot\}\sqcup\mathbb N$ have addition $\max$ and
multiplication given by ordinary addition on $\mathbb N$, with
absorbing zero $\bot$ and multiplicative unit $0\in\mathbb N$.
For a nonzero unital ring $R$, put $F(\bot)=0$ and $F(n)=R$ for
all $n\in\mathbb N$, with identity transitions between nonbottom
fibers and ordinary coefficient products. These operations satisfy the
pair axioms directly. Every nonbottom support is above the unit, so
the ideal of Theorem~\ref{thm:upper-unit-quotient} is zero.

Here $U(Q)=\{0\}$. At support $0$ the unit test gives
$K^{\mathrm{id}}_0=0$. For $n>0$, a context with nonbottom exterior
supports has support $p+n+r>0$ and cannot reach an upper idempotent.
Contexts with a bottom exterior support have zero coefficients.
Therefore $K^{\mathrm{id}}_n=R$ for every $n>0$. This proves strict
containment. The coefficient $1_R$ at support $1\in\mathbb N$ has
nonzero powers at all positive supports, so the global ideal is not
nilpotent. Nevertheless its kernel in each fixed idempotent matrix
fiber has the nilpotence bound just proved.
\end{example}
